
\documentclass[11pt]{article}
\usepackage[T1]{fontenc}
\usepackage[utf8]{inputenc}
\usepackage[english]{babel}
\usepackage{lmodern}
\usepackage{microtype}
\usepackage{amsmath,amssymb,amsthm,mathtools}
\usepackage{mathrsfs}
\usepackage{enumitem}
\usepackage{geometry}
\geometry{margin=1.02in}
\usepackage{hyperref}
\hypersetup{colorlinks=true,linkcolor=blue,citecolor=blue,urlcolor=blue,hypertexnames=false}

\newcommand{\U}{\mathcal{U}}
\newcommand{\V}{\mathcal{V}}
\newcommand{\F}{\mathcal{F}}
\newcommand{\G}{\mathcal{G}}
\newcommand{\Oo}{\mathcal{O}}
\newcommand{\A}{\mathcal{A}}
\newcommand{\M}{\mathcal{M}}
\newcommand{\R}{\mathcal{R}}
\newcommand{\Ff}{\mathcal{F}}
\newcommand{\Gg}{\mathcal{G}}
\newcommand{\Sch}{\operatorname{Sch}}
\newcommand{\Et}{\operatorname{Et}}
\newcommand{\Zar}{\operatorname{Zar}}
\newcommand{\fpqc}{\mathrm{fpqc}}
\newcommand{\fppf}{\mathrm{fppf}}
\newcommand{\Ens}{\mathbf{Set}}
\newcommand{\Ab}{\mathbf{Ab}}
\newcommand{\Cat}{\mathbf{Cat}}
\newcommand{\Top}{\mathbf{Top}}
\newcommand{\Spec}{\operatorname{Spec}}
\newcommand{\Hom}{\operatorname{Hom}}
\newcommand{\Ob}{\operatorname{Ob}}
\newcommand{\Ext}{\operatorname{Ext}}
\newcommand{\Sh}{\operatorname{Sh}}
\newcommand{\esp}{\operatorname{sp}}
\newcommand{\EE}{\mathcal{E}}
\newcommand{\ZZ}{\mathbb{Z}}
\newcommand{\fprod}[3]{{#1}\times_{#3}{#2}}
\newcommand{\tw}[1]{\widetilde{#1}}
\newcommand{\ie}{i.e.\ }
\newcommand{\cf}{cf.\ }
\newcommand{\resp}{respectively}
\newcommand{\SGA}{SGA}
\newcommand{\EGA}{EGA}
\newcommand{\loccit}{loc. cit.}
\newcommand{\qedtext}{\hfill\(\square\)}

\newenvironment{statement}[1]{\par\medskip\noindent\textbf{#1.}\ }{\par\medskip}
\newenvironment{remarktext}[1]{\par\medskip\noindent\textbf{#1.}\ }{\par\medskip}
\newenvironment{prooftext}{\par\noindent\emph{Proof.}\ }{\par\medskip}
\setlist[enumerate]{leftmargin=2.25em}
\setlist[itemize]{leftmargin=2.25em}

\begin{document}

% SGA TRANSLATION PROJECT
% Volume: SGA 4
% Expose: VII
% Range translated: complete expose, source lines 34--1032 of 07/07.tex.
% Status: draft translation, compile-checked, not proofed line-by-line against scan.

\begin{center}
{\Large \textbf{SGA 4, Expos\'e VII}}\\[0.35em]
{\large \textbf{\'Etale Site and \'Etale Topos of a Scheme}}\\[0.25em]
{\normalsize A. Grothendieck}
\end{center}

\bigskip

In the present expos\'e and the following one we develop the most elementary properties relative to the \'etale topology and \'etale cohomology. The developments of the present expos\'e concern certain properties which, in essence, remain valid for other quite different topologies, such as the fppf topology. In the following expos\'e we shall develop properties rather special to the \'etale topology, depending on the very particular nature of \'etale morphisms.

We partially follow here three oral lectures of J. E. Roos, which he was unable to write up, especially in the proof of VIII, 6.3.

Unless expressly stated otherwise, it will be understood that the schemes considered in the present expos\'e and the following ones are elements of the fixed universe \(\U\).

\section*{1. The \'etale topology}

\subsection*{1.1}
We denote by \((\Sch)\) the category of schemes, as elements of the fixed universe \(\U\). Recall that a morphism \(f:X\to Y\) in \((\Sch)\) is called \emph{\'etale} if it is locally of finite presentation, flat, and if for every \(y\in Y\) the fibre \(X_y\) is discrete and its local rings are finite \emph{separable} extensions of \(k(y)\). Equivalently, \(f\) is locally of finite presentation and, for every \(Y\)-scheme \(Y'\) which is affine in the absolute sense and every closed subscheme \(Y'_0\) defined by a nilpotent ideal, the map
\[
\Hom_Y(Y',X) \longrightarrow \Hom_Y(Y'_0,X)
\]
is bijective. For the principal properties of this notion we refer to \EGA{} IV, §§17 and 18, and provisionally to \SGA{} 1, I and IV.

\subsection*{1.2}
The \emph{\'etale topology} on \((\Sch)\) is the topology generated by the pretopology for which, for each \(X\in (\Sch)\), \(\operatorname{Cov} X\) consists of the families
\[
(X_i \xrightarrow{u_i} X)_{i\in I}
\]
with \(I\in \U\), such that the \(u_i\) are \'etale and \(X=\bigcup_i u_i(X_i)\) set-theoretically.

It is convenient to associate to every \(X\) the full subcategory \(\Et_{/X}\) of \((\Sch)_{/X}\) consisting of the arrows \(X'\to X\) which are \'etale. We endow it with the topology induced, in the sense of Expos\'e III, by the \'etale topology of \((\Sch)\); this is also called the \'etale topology on \(X\). We denote the resulting site by \(X_{\mathrm{\'et}}\), the \'etale site of \(X\). Strictly speaking, it would be preferable to reserve the notation \(X_{\mathrm{\'et}}\) for the topos \(\widetilde{X_{\mathrm{\'et}}}\) defined by this site; for practical reasons we keep here the notation used in the original seminar.

Every morphism in \(\Et_{/X}\) is \'etale. It follows that a family \((X'_i\xrightarrow{u_i}X')_{i\in I}\) in \(X_{\mathrm{\'et}}\) is covering if and only if it is surjective, \ie
\[
\bigcup_i u_i(X'_i)=X'.
\]
As follows immediately from 3.1, \(X_{\mathrm{\'et}}\) is a \(\U\)-site, so that the results of Expos\'es I--VI apply. The topos \(\widetilde{X_{\mathrm{\'et}}}\) of \(\U\)-sheaves on \(X_{\mathrm{\'et}}\) is the \emph{\'etale topos} of \(X\).

\subsection*{1.3}
In what follows, unless expressly stated otherwise, all topological notions considered in \(\Et_{/X}\) or in \((\Sch)\) are understood with respect to the \'etale topology. Moreover, in language and notation we shall often write \(X\) instead of \(X_{\mathrm{\'et}}\) or \(\Et_{/X}\). Thus a ``sheaf on \(X\)'' will mean a sheaf on the site \(X_{\mathrm{\'et}}\). We denote the category of these sheaves by \(\widetilde{X_{\mathrm{\'et}}}\); it is a \(\U\)-topos. If \(F\) is an abelian sheaf on \(X\), its cohomology groups will be denoted \(H^i(X,F)\), which in Expos\'e V would be written \(H^i(X_{\mathrm{\'et}},F)\).

\subsection*{1.4}
Let \(f:X\to Y\) be a morphism of schemes. The inverse-image functor
\[
Y' \longmapsto Y'\times_Y X
\]
induces a functor
\[
f^*_{\mathrm{\'et}}:\Et_{/Y}\longrightarrow \Et_{/X},
\]
which commutes with finite projective limits and sends covering families to covering families. Hence it is a morphism of sites
\[
f^*_{\mathrm{\'et}}:X_{\mathrm{\'et}}\longrightarrow Y_{\mathrm{\'et}},
\]
and therefore gives a functor on categories of sheaves
\[
f^{\mathrm{\'et}}_*:\widetilde{X_{\mathrm{\'et}}}\longrightarrow \widetilde{Y_{\mathrm{\'et}}},
\qquad
f^{\mathrm{\'et}}_*(F)=F\circ f^*_{\mathrm{\'et}}.
\]
Moreover \(f^{\mathrm{\'et}}_*\) admits a left adjoint
\[
f^*_{\mathrm{\'et}}:\widetilde{Y_{\mathrm{\'et}}}\longrightarrow \widetilde{X_{\mathrm{\'et}}}
\]
extending the functor just defined on sites, and commuting with arbitrary inductive limits and finite projective limits. In other words \(f^*_{\mathrm{\'et}}\) defines a morphism of topoi
\[
f_{\mathrm{\'et}}:\widetilde{X_{\mathrm{\'et}}}\longrightarrow \widetilde{Y_{\mathrm{\'et}}}.
\]

For a composite \(X\xrightarrow{f}Y\xrightarrow{g}Z\), there are canonical isomorphisms
\[
(gf)^{\mathrm{\'et}}_*\simeq g^{\mathrm{\'et}}_*f^{\mathrm{\'et}}_*,
\qquad
(gf)^*_{\mathrm{\'et}}\simeq f^*_{\mathrm{\'et}}g^*_{\mathrm{\'et}},
\qquad
(gf)_{\mathrm{\'et}}\simeq g_{\mathrm{\'et}}f_{\mathrm{\'et}}.
\]
Thus one obtains a pseudo-functor from \((\Sch)\) to the category of topoi in a universe \(\U'\) containing \(\U\).

\subsection*{1.5. Notation}
We shall often write \(f^*\) and \(f_*\) instead of \(f^*_{\mathrm{\'et}}\) and \(f^{\mathrm{\'et}}_*\). If \(f:X\to Y\) is a morphism of schemes, the functors \(R^i f^{\mathrm{\'et}}_*\) will therefore be denoted simply by \(R^i f_*\). With this notation the Leray spectral sequence for \(f\) and an abelian sheaf \(F\) on \(X\) is
\[
H^*(X,F) \Longleftarrow E_2^{p,q}=H^p(Y,R^qf_*(F)).
\]
Likewise, for two morphisms \(f:X\to Y\) and \(g:Y\to Z\), one has the Leray spectral sequence for the composite functor,
\[
R^*(gf)_*(F) \Longleftarrow E_2^{p,q}=R^pg_*(R^qf_*(F)).
\]

\subsection*{1.6}
When \(f:X\to Y\) is itself \'etale, \(X\) may be regarded as an object of \(Y_{\mathrm{\'et}}\), and there is a canonical isomorphism of sites
\[
X_{\mathrm{\'et}} \xrightarrow{\sim} (Y_{\mathrm{\'et}})_{/X}.
\]
Under this isomorphism, the functor \(f^*_{\mathrm{\'et}}:Y'\mapsto Y'\times_Y X\) identifies with the internal base-change functor in \(\Et_{/Y}\). Consequently
\[
f^*:\widetilde{Y_{\mathrm{\'et}}}\longrightarrow \widetilde{X_{\mathrm{\'et}}}
\]
is isomorphic to the restriction functor
\[
f^*(F)\simeq F\circ i_{X/Y},
\]
where \(i_{X/Y}:X_{\mathrm{\'et}}\to Y_{\mathrm{\'et}}\) is the evident functor which regards an \'etale \(X\)-scheme \(u:X'\to X\) as an \'etale \(Y\)-scheme by means of \(fu:X'\to Y\).

\subsection*{1.7. Universes}
If \(X\ne\varnothing\), then \(X_{\mathrm{\'et}}\) is not an element of the chosen universe \(\U\). Nevertheless, as already noted, one can find a full subcategory \(C\) of \(X_{\mathrm{\'et}}\), itself an element of \(\U\), satisfying the hypotheses of the comparison lemma of Expos\'e III, so that restriction induces an equivalence
\[
\widetilde{X_{\mathrm{\'et}}}\simeq \widetilde{C}.
\]
Thus the \'etale topos is equivalent to a topos defined by a site \(C\in\U\). Equivalently, \(X_{\mathrm{\'et}}\) is a \(U\)-site, and the results of Expos\'es I--VI apply to it.

\section*{2. Examples of sheaves}

\begin{enumerate}[label=\alph*)]
\item Let \(F\in\Ob((\Sch)_{/X})\) be a scheme over \(X\), and for every \(X'\) over \(X\) put
\[
F(X')=\Hom_X(X',F).
\]
The resulting functor \((\Sch)_{/X}^{\circ}\to\Ens\) is a sheaf for the \'etale topology, and even for the finer fpqc topology studied in \SGA{} 3, IV. In other words, the \'etale topology on \((\Sch)\) is less fine than the canonical topology; this is essentially the content of \SGA{} 1, VIII, 5.1, and also of \SGA{} 3, IV, 6.3.1. A fortiori its restriction to \(X_{\mathrm{\'et}}\) is a sheaf; it will again be denoted by \(F\) when no confusion can arise.

One must remember that, if \(X\ne\varnothing\), the functor \(F\mapsto \varphi(F)\), from schemes over \(X\) to sheaves on \(X_{\mathrm{\'et}}\), is neither fully faithful nor even faithful; one cannot reconstruct \(F\), up to isomorphism, from the sheaf \(\varphi(F)\). For instance, when \(X=\Spec k\) with \(k\) algebraically closed, the sheaf remembers only the underlying set of \(F\). Compare Expos\'e IV.

The functor
\[
( \Sch)_{/X}\longrightarrow \widetilde{X_{\mathrm{\'et}}}
\]
so obtained commutes with finite projective limits. Hence, if \(F\) is a group scheme, or an object with similar finite-limit structure, the associated sheaf is a sheaf of groups, and so on. Its restriction
\[
\Et_{/X}\longrightarrow \widetilde{X_{\mathrm{\'et}}}
\]
is just the canonical Yoneda functor; it identifies \(\Et_{/X}\) with a full subcategory of the \'etale topos. We generally identify an object of \(X_{\mathrm{\'et}}\) with the corresponding sheaf, denoted \(\widetilde{X'}\) or simply \(X'\).

Evidently
\[
H^0(X,F)=\Hom_{(\Sch)_{/X}}(X,F).
\]
If \(F\) is a group prescheme over \(X\), commutative if necessary so that \(H^1(X,F)\) is within the framework of Expos\'e V, the usual arguments show that \(H^1(X,F)\) is canonically isomorphic to the set of isomorphism classes of principal homogeneous sheaves of sets under \(F\), or \(F\)-torsors, on \(X_{\mathrm{\'et}}\). If \(F\) is affine over \(X\), \SGA{} 1, VIII, 2.1, together with standard arguments, shows that \(H^1(X,F)\) is also the set of classes of schemes \(P\) over \(X\), endowed with a right action of \(F\), which are principal homogeneous bundles over \(X\) for the \'etale topology, equivalently representable \(F\)-torsors.
\end{enumerate}

\begin{remarktext}{Remarks 2.1}
Let \(\varphi_X(Z)\) denote the sheaf on \(X_{\mathrm{\'et}}\) associated with a scheme \(Z\) over \(X\), and let \(f:Y\to X\) be a morphism of schemes. There is an evident homomorphism, functorial in \(Z\),
\[
f^*\varphi_X(Z)\longrightarrow \varphi_Y(Z_Y),
\qquad Z_Y=Z\times_XY,
\]
or, equivalently, an evident homomorphism
\[
\varphi_X(Z)\longrightarrow f_*\varphi_Y(Z_Y),
\]
given on \(X'\in\Ob X_{\mathrm{\'et}}\) by
\[
\Hom_X(X',Z)\longrightarrow \Hom_Y(X'_Y,Z_Y).
\]
In general this homomorphism is not an isomorphism: formation of the \'etale sheaf associated with a relative scheme does not commute with inverse-image functors. It is, however, an isomorphism when \(Z\) is \'etale over \(X\). Likewise \(\varphi_X:\Sch_{/X}\to\widetilde{X_{\mathrm{\'et}}}\) is not faithful when \(X\ne\varnothing\), though its restriction to \(X_{\mathrm{\'et}}\) is fully faithful, since the topology of \(X_{\mathrm{\'et}}\) is less fine than its canonical topology.

\begin{enumerate}[label=\alph*)]
\setcounter{enumi}{1}
\item The preceding functor also commutes with direct sums indexed by sets in \(\U\). In particular, if \(I_X\) denotes the constant \(X\)-scheme defined by a set \(I\), the associated sheaf is just the constant sheaf \(I_{X_{\mathrm{\'et}}}\). Since \(I\mapsto I_X\) also commutes with finite projective limits, it carries groups to groups, commutative groups to commutative groups, and so on. If \(G\) is an ordinary commutative group, we simply write \(H^i(X,G)\) for \(H^i(X,G_X)\).

If, for example, \(G\) is finite, then \(G_X\) is finite and hence affine over \(X\), and the final observation in (a) gives an interpretation of \(H^1(X,G)\) as the set of classes of principal Galois coverings with group \(G\). When \(X\) is connected and endowed with a geometric point \(a\), the fundamental group \(\pi_1(X,a)\) gives the canonical isomorphism, for \(G\) commutative,
\[
H^1(X,G)\simeq \Hom(\pi_1(X,a),G).
\]
One may say that, on passing from Zariski cohomology to the \'etale topology, one has done what is necessary to obtain the correct \(H^1\) for a finite constant coefficient group. It is a remarkable fact, proved later in this seminar, that this is also enough to obtain the correct \(H^i(X,G)\) for every torsion coefficient group, at least when \(G\) is prime to the residue characteristics of \(X\).

\item Let \(F\) be an \(\Oo_X\)-module on \(X\) for the Zariski topology. With the notation of \SGA{} 3, I, 4.6, define
\[
W(F):(\Sch)_{/X}^{\circ}\longrightarrow \Ens
\]
by
\[
W(F)(X')=\Gamma(X',F\otimes_{\Oo_X}\Oo_{X'}).
\]
This is again a sheaf for the \'etale topology, and even for the fpqc topology, by \SGA{} 1, VIII, 1.6 and \SGA{} 3, IV, 6.3.1. Its restriction to \(X_{\mathrm{\'et}}\), again denoted \(W(F)\), satisfies
\[
H^0(X,W(F))=\Gamma(X,F)=H^0(X,F).
\]
If \(F\) is quasi-coherent, more is true; see 4.3.
\end{enumerate}
\end{remarktext}

\section*{3. Generators of the \'etale topos. Cohomology of an inductive limit of sheaves}

\begin{statement}{Proposition 3.1}
Let \(X\) be a scheme and let \(C\) be a full subcategory of \(X_{\mathrm{\'et}}\) such that, for every affine \'etale \(X\)-scheme \(X'\), \(C\) contains an object isomorphic to \(X'\). Then \(C\) is a topological generating family of the site \(X_{\mathrm{\'et}}\), and hence a generating family of the topos \(\widetilde{X_{\mathrm{\'et}}}\). Moreover restriction
\[
\widetilde{X_{\mathrm{\'et}}}\longrightarrow \widetilde{C}
\]
is an equivalence of categories, where \(C\) is endowed with the topology induced from \(X_{\mathrm{\'et}}\).
\end{statement}

This is immediate from the comparison lemma.

\begin{statement}{Corollary 3.2}
Assume that \(X\) is quasi-separated. Let the restricted \'etale site of \(X\) be the full subcategory \(C\) of \(X_{\mathrm{\'et}}\) consisting of schemes \'etale and of finite presentation over \(X\), with the topology induced from \(X_{\mathrm{\'et}}\). Then:
\begin{enumerate}[label=(\roman*)]
\item \(C\) is stable under fibre products, and is a site of finite type if \(X\) is quasi-compact;
\item restriction \(\widetilde{X_{\mathrm{\'et}}}\to\widetilde{C}\) is an equivalence of categories.
\end{enumerate}
\end{statement}

Assertion (i) is trivial. Assertion (ii) follows from 3.1, since quasi-separatedness of \(X\) implies that an \'etale \(X\)-scheme \(X'\) which is quasi-compact, for example affine, is quasi-compact over \(X\); if it is also locally of finite type over \(X\), it is of finite presentation over \(X\).

\begin{statement}{Proposition 3.3}
Assume that \(X\) is quasi-compact and quasi-separated. Then the functors \(H^q(X_{\mathrm{\'et}},F)\), on the category of abelian sheaves on \(X_{\mathrm{\'et}}\), commute with filtered inductive limits.
\end{statement}

Indeed, by 3.2(ii) one may replace \(X_{\mathrm{\'et}}\) by the restricted site \(C\). Since \(X\) is quasi-compact, \(X\in\Ob C\), and the conclusion follows from 3.2(i) and VI, 6.1.2(3).

\section*{4. Comparison with other topologies}

\subsection*{4.0}
First note that the examples of sheaves on \(X_{\mathrm{\'et}}\) considered in §2 are naturally restrictions of sheaves defined on \((\Sch)_{/X}\), endowed with its \'etale topology, or even with the fpqc topology. In general let
\[
u^*:X_{\mathrm{\'et}}\longrightarrow (\Sch)_{/X}
\]
be the inclusion functor. It is continuous and commutes with finite projective limits, and hence defines a morphism of sites
\[
u:(\Sch)_{/X}\longrightarrow X_{\mathrm{\'et}},
\]
with the corresponding direct-image functor
\[
u_*:\widetilde{(\Sch)_{/X}}\longrightarrow \widetilde{X_{\mathrm{\'et}}},
\qquad
u_*(F)=F\circ u,
\]
and a left adjoint
\[
u^*:\widetilde{X_{\mathrm{\'et}}}\longrightarrow \widetilde{(\Sch)_{/X}}.
\]
In complete rigor one must take account of universes here, since \((\Sch)_{/X}\) is a \(\U\)-category but not a \(\U\)-site. One may either work in a larger universe or replace \((\Sch)_{/X}\) by a suitable full \(U\)-small subcategory stable under finite projective limits and containing a comparison subcategory as in 3.1.

\begin{statement}{Proposition 4.1}
\begin{enumerate}[label=(\roman*)]
\item The functor \(u^*\) is fully faithful. Thus, for every sheaf \(G\) on \(X_{\mathrm{\'et}}\), the canonical homomorphism \(G\to u_*u^*G\) is an isomorphism.
\item If \(G\) is an abelian sheaf on \(X_{\mathrm{\'et}}\) and \(F\) is an abelian sheaf on the big \'etale site \((\Sch)_{/X}\), then the canonical comparison maps
\[
H^*(X_{\mathrm{\'et}},G)\xrightarrow{\sim}H^*((\Sch)_{/X},u^*G),
\qquad
H^*(X_{\mathrm{\'et}},u_*F)\xrightarrow{\sim}H^*((\Sch)_{/X},F)
\]
are isomorphisms.
\end{enumerate}
\end{statement}

Indeed, the functor \(u\) is a comorphism, and one applies Expos\'e III. Thus, for the study of sheaf cohomology, it is essentially equivalent to work with the small \'etale site \(X_{\mathrm{\'et}}\) or with the big \'etale site \((\Sch)_{/X}\).

\subsection*{4.2}
It is also useful to introduce on \((\Sch)\), and hence on each \((\Sch)_{/X}\), topologies other than the \'etale topology. The most useful among these are those defined in \SGA{} 3, IV, 6.3. The coarsest is the \emph{Zariski topology}, whose covering families are the surjective families of open immersions; it is less fine than the \'etale topology. The finest of these is the faithfully flat quasi-compact topology, abbreviated fpqc, the least fine topology for which Zariski coverings and faithfully flat quasi-compact morphisms are covering; it is finer than the \'etale topology. As already noted, the examples of sheaves on \((\Sch)_{/X}\) considered in §2 are already sheaves for the fpqc topology.

\subsubsection*{4.2.1}
One must beware, however, that if \(F\) is an abelian sheaf on \((\Sch)_{/X}\) for the fpqc topology, or for a topology such as fppf considered in \SGA{} 3, the cohomology groups of \(F\) for the fpqc topology are not always isomorphic to those for the \'etale topology. This may already occur when \(X=\Spec k\) and \(F\) is represented by an algebraic group over \(k\), even for \(H^1\). In general the \'etale topology gives the correct cohomology groups for coefficient groups which are \'etale, or more generally smooth, group schemes over \(X\); but this is no longer true for group schemes such as radicial groups, for which the \'etale topology is too coarse and must be replaced by the fpqc or fppf topology.

\subsubsection*{4.2.2}
As an example of the relation between cohomology for different topologies, consider the Zariski and \'etale topologies. Let \(X_{\Zar}\) be the site of Zariski opens of \(X\). There is a canonical inclusion functor
\[
u^*:X_{\Zar}\longrightarrow X_{\mathrm{\'et}},
\]
defining a morphism of sites
\[
u:X_{\mathrm{\'et}}\longrightarrow X_{\Zar},
\]
and hence direct and inverse image functors
\[
u_*:\widetilde{X_{\mathrm{\'et}}}\longrightarrow \widetilde{X_{\Zar}},
\qquad
u^*:\widetilde{X_{\Zar}}\longrightarrow \widetilde{X_{\mathrm{\'et}}}.
\]
Geometrically, the pair \((u_*,u^*)\) is a morphism of topoi
\[
u:\widetilde{X_{\mathrm{\'et}}}\longrightarrow \widetilde{X_{\Zar}}.
\]
It yields a morphism of cohomological functors
\[
H^*(X_{\Zar},u_*F)\longrightarrow H^*(X_{\mathrm{\'et}},F)
\]
and the Leray spectral sequence
\[
H^*(X_{\mathrm{\'et}},F)\Longleftarrow H^p(X_{\Zar},R^qu_*F),
\]
for every abelian sheaf \(F\) on \(X_{\mathrm{\'et}}\). This spectral sequence summarizes the general relation between \'etale and Zariski cohomology. For coefficient sheaves such as constant sheaves it is generally far from trivial; indeed in general \(R^q u_*F\ne 0\), already over the spectrum of a field. However:

\begin{statement}{Proposition 4.3}
Let \(F\) be a sheaf of modules on the prescheme \(X\) for the Zariski topology, and let \(F_{\mathrm{\'et}}\) be the corresponding sheaf on \(X_{\mathrm{\'et}}\). There is a natural morphism of cohomological functors
\[
H^*(X,F)\longrightarrow H^*(X_{\mathrm{\'et}},F_{\mathrm{\'et}}).
\]
If \(F\) is quasi-coherent, this morphism is an isomorphism.
\end{statement}

The left hand side is \(H^*(X_{\Zar},u_*F_{\mathrm{\'et}})\). It is enough to show
\[
R^q u_*(F_{\mathrm{\'et}})=0\quad (q>0).
\]
Using the standard computation of \(R^qu_*\) and the fact that affine opens form a base of the topology of \(X\), this follows from the following corollary.

\begin{statement}{Corollary 4.4}
If \(X\) is affine and \(F\) is quasi-coherent, then
\[
H^q(X_{\mathrm{\'et}},F_{\mathrm{\'et}})=0\quad (q>0).
\]
\end{statement}

Let \(C=X_{\mathrm{\'et}}\), and let \(C'\) be the full subcategory of affine objects. By 3.1,
\[
H^q(C,F_{\mathrm{\'et}})\simeq H^q(C',F|_{C'}).
\]
Thus it suffices to prove that \(F|_{C'}\) is \(C'\)-acyclic, or equivalently that it satisfies the criterion of Expos\'e V, 4.3: for every \(X'\in\Ob C'\) and every covering family \(\R=(X'_i\to X')_i\) in \(C'\), one has \(H^q(\R,F)=0\) for \(q>0\). We may suppose the family finite, and then replace the \(X'_i\) by their sum; hence the covering consists of a single surjective \'etale morphism \(X''\to X'\) of affine schemes. We are reduced to proving that, if \(f:X\to Y\) is a surjective \'etale morphism of affine schemes and \(F\) is a quasi-coherent module on \(X\), then
\[
H^q(X/Y,F)=0\quad(q>0).
\]
This is proved in the theory of faithfully flat descent.

\begin{remarktext}{Remarks 4.5}
In fact, for the last vanishing it is enough that \(X\to Y\) be surjective and flat, not necessarily \'etale. The preceding proof therefore also gives analogous isomorphisms with the \'etale site replaced by \((\Sch)_{/X}\) endowed with any of the topologies finer than the Zariski topology considered in \SGA{} 3, IV, 6.3, for example the fpqc topology.
\end{remarktext}

\section*{5. Cohomology of a projective limit of schemes}

\subsection*{5.1}
Let \(I\) be an increasing filtered preordered set, and let
\[
\mathfrak{X}=(X_i)_{i\in I}
\]
be a projective system of schemes, with affine transition morphisms \(u_{ij}:X_i\to X_j\) for \(i\ge j\). Recall from \EGA{} IV, 8, that under these hypotheses the projective limit
\[
X=\varprojlim_i X_i
\]
exists in the category of schemes, indeed in the category of locally ringed spaces. It may be constructed as follows. Choose \(i_0\in I\). For \(i\ge i_0\), write over \(X_{i_0}\)
\[
X_i=\Spec \A_i,
\]
where \(\A_i\) is a quasi-coherent algebra; the \(\A_i\) form an inductive system of such algebras. With
\[
\A=\varinjlim_i \A_i,
\]
one may take
\[
X=\Spec \A.
\]
Moreover, if \(\esp S\) denotes the underlying topological space of a scheme \(S\), then the canonical map
\[
\esp X\longrightarrow \varprojlim_i \esp(X_i)
\]
is a homeomorphism, and the canonical morphism of sheaves of rings
\[
\varinjlim_i \gamma_i^{-1}\Oo_{X_i}\longrightarrow \Oo_X
\]
is bijective, where \(\gamma_i:X\to X_i\) is the canonical continuous map.

\subsection*{5.2}
Informally, the content of \EGA{} IV, 8 and 9 may be summarized as follows: when \(X_{i_0}\) is quasi-compact and quasi-separated, every schematic datum on \(X\) which is of finite presentation over \(X\) is equivalent to a datum of the same nature over some \(X_i\), for sufficiently large \(i\). Thus one proves:

\begin{enumerate}[label=(\alph*)]
\item If \(i_1\in I\) and \(X'_{i_1}\), \(X''_{i_1}\) are two schemes of finite presentation over \(X_{i_1}\), put, for \(i\ge i_1\),
\[
X'_i=X'_{i_1}\times_{X_{i_1}}X_i,
\qquad
X'=X'_{i_1}\times_{X_{i_1}}X,
\]
and define \(X''_i\) and \(X''\) similarly. Then the canonical map
\[
\varinjlim_{i\ge i_1}\Hom_{X_i}(X'_i,X''_i)\longrightarrow \Hom_X(X',X'')
\]
is bijective.
\item For every scheme \(X'\) of finite presentation over \(X\), there exist an index \(i_1\in I\), a scheme \(X'_{i_1}\) of finite presentation over \(X_{i_1}\), and an \(X\)-isomorphism
\[
X'\simeq X'_{i_1}\times_{X_{i_1}}X.
\]
\end{enumerate}

\subsection*{5.3}
The preceding result, and analogous results of \EGA{} IV, 8, can be expressed in the language of inductive limits of fibred categories introduced in Expos\'e VI. Let \(\Ff\) be the category of morphisms \(f:T\to S\) of finite presentation of schemes, and let
\[
\pi:\Ff\longrightarrow (\Sch)
\]
be the target functor. This is a fibred functor, and its fibres are equivalent to \(\U\)-small categories. Pulling back the fibred category \(\Ff/(\Sch)\) along the functor
\[
i\longmapsto X_i:\operatorname{cat}(I)\longrightarrow(\Sch)
\]
gives a fibred category
\[
\pi_{\mathfrak X}:\Ff_{\mathfrak X}\longrightarrow \operatorname{cat}(I),
\]
whose fibre at \(i\) is the full subcategory \(\Ff_{X_i}\) of \((\Sch)_{/X_i}\) consisting of schemes of finite presentation over \(X_i\); the change-of-base functor from \(i\) to \(j\) is the usual base change \(X'_i\mapsto X'_i\times_{X_i}X_j\). Associating to \(X'_i\) its inverse image over \(X\),
\[
\varphi(X'_i)=X'_i\times_{X_i}X,
\]
we get a natural functor \(\varphi:\Ff_{\mathfrak X}\to\Ff_X\) which sends cartesian morphisms to isomorphisms. Hence it factors uniquely through a canonical functor
\[
\psi:\varinjlim_{\Ff_{\mathfrak X}/\operatorname{cat}(I)} \Ff_{\mathfrak X}
\longrightarrow \Ff_X.
\]
The result of \EGA{} IV, 8 recalled above says precisely that \(\psi\) is an equivalence of categories.

\subsection*{5.4}
Statement (a) of 5.2 can be refined in various ways by introducing a class \(\M\) of morphisms of schemes, stable under base change, and asserting that for a given \(X_{i_1}\)-morphism \(u_{i_1}:X'_{i_1}\to X''_{i_1}\), the corresponding morphism \(u':X'\to X''\) has property \(\M\) if and only if, for some \(i\ge i_1\), the morphism \(u_i:X'_i\to X''_i\) has property \(\M\). This holds, for example, when \(\M\) is one of the following classes: proper, projective, quasi-projective, affine, finite, quasi-finite, radicial, surjective, flat, smooth, \'etale, or unramified morphisms. The corresponding statements are found in \EGA{} IV, §8 for the first cases, §11 for flatness, and §17 for the last three cases.

\subsection*{5.5}
Replace \(\Ff\to(\Sch)\) by the fibred subcategory \(\Gg\) consisting of \'etale morphisms of finite presentation. Its fibre \(\Gg_S\) over a scheme \(S\) is the category of schemes \'etale and of finite presentation over \(S\). We form similarly the fibred category \(\Gg_{\mathfrak X}\) over \(\operatorname{cat}(I)\), whose fibre at \(i\) is \(\Gg_{X_i}\). By 3.2(ii), the \'etale topoi \(\widetilde{(X_i)_{\mathrm{\'et}}}\) and \(\widetilde{X_{\mathrm{\'et}}}\) are canonically equivalent to \(\widetilde{\Gg_{X_i}}\) and \(\widetilde{\Gg_X}\), where each \(\Gg_S\) carries the topology induced from \(S_{\mathrm{\'et}}\). By 3.2(i), each \(\Gg_{X_i}\) is a site of finite type, and equivalent to a \(\U\)-small site. With these topologies, \(\Gg\) becomes a fibred site over \(\operatorname{cat}(I)\).

\begin{statement}{Lemma 5.6}
The canonical functor
\[
\varinjlim_{\Gg_{\mathfrak X}/\operatorname{cat}(I)}\Gg_{\mathfrak X}
\longrightarrow \Gg_X
\]
defines an equivalence, respecting the topologies, from the inductive-limit site of the restricted \'etale sites of the \(X_i\) to the restricted \'etale site of \(X=\varprojlim_i X_i\).
\end{statement}

The categorical equivalence is the assertion recalled in 5.4 for \'etale morphisms. The assertion about topologies follows in the same way by taking \(\M\) to be the class of surjective morphisms of schemes.

\begin{statement}{Theorem 5.7}
Let \(\mathfrak X=(X_i)_{i\in I}\) be a projective system of schemes indexed by an increasing filtered preordered set. Assume that the transition morphisms \(u_{ij}:X_j\to X_i\) are affine, so that \(X=\varprojlim X_i\) is defined as in 5.1, and that the \(X_i\) are quasi-compact and quasi-separated. Let \(F\) be an abelian sheaf on the fibred site \(\Gg\), \ie a functor \(\Gg^\circ\to\Ab\) whose restriction \(F_i\) to each \(\Gg_i\) is a sheaf on \(X_i\). Let
\[
F_\infty=\varinjlim_i u_i^*(F_i)
\]
be the sheaf induced on \(X\), where \(u_i:X\to X_i\) is the canonical morphism. Then the canonical homomorphisms
\[
\varinjlim_i H^n(X_i,F_i)\longrightarrow H^n(X,F_\infty)
\]
are isomorphisms.
\end{statement}

In view of 5.6, this is the conjunction of VI, 8.7.3 and VI, 8.5.2. We shall mainly use the following particular case.

\begin{statement}{Corollary 5.8}
With the preceding hypotheses and notation for \((X_i)_{i\in I}\), \(X\), \(u_{ij}\), and \(u_i\), let \(i_0\in I\), and let \(F_{i_0}\) be an abelian sheaf on \(X_{i_0}\). For \(i\ge i_0\), let \(F_i=u_{i_0i}^*F_{i_0}\), and let \(F_\infty=u_i^*F_{i_0}\) on \(X\). Then the canonical homomorphisms
\[
\varinjlim_i H^n(X_i,F_i)\longrightarrow H^n(X,F_\infty)
\]
are isomorphisms.
\end{statement}

The following case is often useful and belongs to 5.7, but not to 5.8.

\begin{statement}{Corollary 5.9}
With the hypotheses and notation of 5.7, let \(G_{i_0}\) be a commutative group scheme locally of finite presentation over \(X_{i_0}\). For \(i\ge i_0\), put \(G_i=G_{i_0}\times_{X_{i_0}}X_i\), and likewise \(G_\infty=G_{i_0}\times_{X_{i_0}}X\). Then the canonical homomorphisms
\[
\varinjlim_i H^n(X_i,G_i)\longrightarrow H^n(X,G_\infty)
\]
are isomorphisms.
\end{statement}

Indeed, \(G_{i_0}\) defines a sheaf on \((\Sch)_{/X_{i_0}}\). We may suppose that \(i_0\) is an initial object of \(I\). Composition with the canonical functor \(\Gg\to(\Sch)_{/X_{i_0}}\) gives a contravariant functor \(F\) on \(\Gg\), whose restriction to \(\Gg_i\) is canonically the sheaf on \(X_i\) defined by \(G_i\). Hence \(F\) is a sheaf on \(\Gg\), and 5.7 applies. The hypothesis that \(G_{i_0}\) is locally of finite presentation over \(X_{i_0}\) ensures that the resulting \(F_\infty\) is the sheaf defined by \(G_\infty\), by \EGA{} IV, 8.8.2(i) and Lemma 5.6.

\begin{statement}{Corollary 5.10}
With the hypotheses and notation of 5.7, let \(F\) be a sheaf of commutative groups on \(X\). Then the canonical homomorphisms
\[
\varinjlim_i H^n(X_i,u_{i*}F)\longrightarrow H^n(X,F)
\]
are isomorphisms.
\end{statement}

The statements 5.7--5.10 generalize, by VI, 8.7.9, to corresponding statements for \(\Ext^i\) of sheaves of modules; we shall not repeat these special forms here. Instead we record, for later reference, variants of 5.8 and 5.9 in terms of the functors \(R^nf_*\).

\begin{statement}{Corollary 5.11}
Let \(I\) be an increasing filtered preordered set. Let \((X_i)_{i\in I}\) and \((Y_i)_{i\in I}\) be two projective systems of schemes with affine transition morphisms, let \(X=\varprojlim_iX_i\), \(Y=\varprojlim_iY_i\), and let \((f_i)_{i\in I}\) be a projective system of quasi-compact and quasi-separated morphisms \(f_i:X_i\to Y_i\). Let \(i_0\in I\), and let \(F_{i_0}\) be a sheaf on \(X_{i_0}\). For \(i\ge i_0\), let \(F_i\) be its inverse image on \(X_i\); similarly let \(f:X\to Y\) be induced by the \(f_i\), and let \(F\) be the inverse image of \(F_{i_0}\) on \(X\). Then the canonical homomorphism
\[
\varinjlim_i v_i^*(R^nf_{i*}F_i)\longrightarrow R^nf_*F
\]
is an isomorphism, where \(v_i:Y\to Y_i\) is the canonical morphism.
\end{statement}

The question is local on \(Y_{i_0}\), and one may suppose \(Y_{i_0}\) affine and \(i_0\) initial. Then all the \(Y_i\) and \(X_i\) are quasi-compact and quasi-separated, and the assertion follows from VI, 8.7.3, reducing to 5.8 and using the computation of \(R^nf_*\) as sheaves associated with presheaves.

\begin{statement}{Corollary 5.12}
The same statement as 5.11 holds when \(F_{i_0}\) is a commutative group scheme locally of finite presentation over \(X_{i_0}\), and \(F_i\), \(F\) are its inverse images in the sense of base change for schemes.
\end{statement}

Thus 5.9, respectively 5.12, is not a special case of 5.8, respectively 5.11; see §2(a) and Remark 2.1.

\begin{statement}{Corollary 5.13}
With the notation of 5.11, let \(F\) be a sheaf of abelian groups on \(X\). Then the canonical homomorphisms
\[
\varinjlim_i v_i^*(R^nf_{i*}(u_{i*}F))\longrightarrow R^nf_*F
\]
are isomorphisms, where \(u_i:X\to X_i\) and \(v_i:Y\to Y_i\) are the canonical morphisms.
\end{statement}

The proof is essentially the same as before.

\begin{remarktext}{Remarks 5.14}
\begin{enumerate}[label=\alph*)]
\item The preceding limit-passage results remain valid for \(H^0\), respectively for direct images \(f_*\), with sheaves of sets rather than abelian sheaves, and for \(H^1\), respectively \(R^1f_*\), with sheaves of not necessarily commutative groups; the non-commutative \(H^1\) and \(R^1f_*\) are defined as usual in terms of torsors or principal homogeneous bundles. This can be checked directly in the general context of Expos\'e VI. It should also be possible, and useful, to give a variant for the non-commutative \(H^2\) of ``links'' studied by Giraud.

\item The limit-passage results for fibred sites developed in Expos\'e VI only assume that the base category is filtered; it is not necessary that it come from a filtered preordered set. An arbitrary filtered index category is also the natural framework for the statements developed in the present section. We used the more restrictive framework only in order to give exact references to \EGA{} IV, where, unfortunately, limit-passage questions from §8 onward are stated with index categories defined by filtered preordered sets. We shall nevertheless use henceforth, without further comment, the corresponding \EGA{} results and the results of the present section for arbitrary filtered index categories, since the proofs in \loccit{} apply essentially without change.
\end{enumerate}
\end{remarktext}

\end{document}
